% Generated draft: requires mathematical and source review.
\documentclass[11pt]{article}
\usepackage[margin=25mm]{geometry}
\usepackage{fontspec}
\setmainfont{Latin Modern Roman}
\usepackage{amsmath,amssymb,mathtools,unicode-math}
\setmathfont{Latin Modern Math}
\usepackage{enumitem,needspace,xurl,hyperref}
\setlength{\parindent}{0pt}
\setlength{\parskip}{5pt}
\setlength{\emergencystretch}{3em}
\newcommand{\N}{\mathbb{N}}
\newcommand{\Z}{\mathbb{Z}}
\newcommand{\Q}{\mathbb{Q}}
\newcommand{\R}{\mathbb{R}}
\newcommand{\C}{\mathbb{C}}
\newcommand{\F}{\mathbb{F}}
\newcommand{\A}{\mathbb{A}}
\newcommand{\PP}{\mathbb{P}}
\newcommand{\E}{\mathbb{E}}
\newcommand{\eps}{\varepsilon}
\newcommand{\dd}{\,\mathrm{d}}
\begin{document}
\section*{UnsolvedMath ID: 30001007}
\textbf{Source problem code:} catalogue entry
\subsection{Definitions and mathematical statement}
Fix integers \(r,g\geq0\), \(d\geq1\), and ramification orders \((e_1,\ldots,e_r)\) satisfying the extracted relation
\[
 \sum_{i=1}^r(e_i-1)=2d-2-g.
\]
In the standard characteristic-zero Hurwitz interpretation, a point of \(\mathcal H\) represents a degree-\(d\) cover \(f:C\to\PP^1\) from a smooth connected genus-\(g\) curve. At the \(i\)-th distinguished branch value there is one ramification point \(p_i\) of index \(e_i\), with all remaining points over that value unramified; equivalently, its local monodromy is one \(e_i\)-cycle with the other letters fixed. There are exactly \(3g\) further simple branch values, each with one index-two ramification point. Distinctness, labelling, and the treatment of index-one distinguished points must be specified by the intended Hurwitz space.

The target \(\PP^1\) and characteristic-zero/tame interpretation are not explicit in the extraction, but they explain its numerical relation: Riemann--Hurwitz gives
\[
 2g-2=-2d+\sum_i(e_i-1)+3g.
\]
This consistency does not prove nonemptiness for arbitrary numerical data. Admissible ramification indices lie between one and \(d\), with index at least two for genuine branch points; the source does not explicitly impose those ranges.

Using the distinguished source ramification points as the \(r\) retained markings, the forgetful map is
\[
 F:\mathcal H\longrightarrow\mathcal M_{g,r},\qquad
 [f:C\to\PP^1;p_1,\ldots,p_r;\text{additional data}]
 \longmapsto[C;p_1,\ldots,p_r].
\]
It forgets the covering map and the \(3g\) simple ramification markings. Retaining the target branch values instead would not define an \(r\)-pointed source curve. The question is whether every irreducible component \(T\) of the intended \(\mathcal H\) is dominant:
\[
 \forall T\in\operatorname{Irr}(\mathcal H),\qquad
 \overline{F(T)}=\mathcal M_{g,r}.
\]
This is stronger than dominance of the union or existence of one dominating component. It makes no surjectivity assertion at every pointed curve. An empty Hurwitz space has no components, so the component assertion alone does not assert existence of covers.

The precise field, Hurwitz equivalence convention, fixed versus varying branch data, and low \((g,r)\) moduli conventions are missing. In particular, the source does not impose the usual stable pointed-curve range \(2g-2+r>0\).

The supplied assessment dated 2026-08-21 reports a conditional higher-genus result of Liu--Osserman requiring at least \(3g\) simple branch points, while treating the full component-dominance question as unresolved. The record's exactly \(3g\) condition is retained; the other hypotheses of the conditional result are not supplied, so the inequality alone does not justify identifying it with a full solution.
\subsection{Short English statement}
For the stated pure-cycle ramification data and exactly \(3g\) additional simple branch points, must every Hurwitz-space component cover a dense set of pointed genus-\(g\) curves after forgetting the map? The retained markings are on the source curve, and a reported conditional theorem is not enough without its additional hypotheses.
\subsection{Sources}
\begin{itemize}
\item Local UnsolvedMath registry, permanent ID \(30001007\), exact source code \texttt{catalogue entry}, proposed year 2008; no individual proposer is supplied.
\item \emph{Arithmetic Algebraic Geometry (2009)}, Oberwolfach Reports, DOI 10.4171/owr/2008/35, pp. 1979--2026. \url{https://doi.org/10.4171/owr/2008/35}.
\item Supplied reference labelled ``pure-cycle Hurwitz-space results'': \url{https://arxiv.org/abs/math/0609118}. The assessment attributes conditional higher-genus progress to Liu--Osserman.
\item The supplied \texttt{corrected\_verified} note reports checking content against the official report. The open label and assessment checked 2026-08-21 are supplied dated curation, not independently verified; exact theorem hypotheses and current status require review.
\end{itemize}
\end{document}
